\section{The Angels in the Calendar}\label{app:calendar}

The feasts of the angels in the three Roman calendars, with the rank each
book prints. The pre-1955 column gives the rank of the Missal of the 1920
typical edition (\emph{MR}~1920: the Vatican printing and Mame's Tours
printing \latin{iuxta typicam} of 1922), with the Roman Breviary printed
at Tours in 1898 (\emph{BR}) and Bute's English Breviary of 1908 where
they differ; the 1962 column, the Missal of 1962 (\emph{MR}~1962); the
postconciliar column, the General Roman Calendar of the Missal of Paul~VI
with the \emph{Ordo lectionum Missae} of 1981 (\emph{OLM}). Page numbers are
those of the books named. The feasts and their texts are expounded in
\S\ref{sec:liturgy}.

\begin{center}
\scriptsize
\setlength{\tabcolsep}{3pt}
\begin{longtable}{@{}>{\raggedright\arraybackslash}p{0.15\linewidth}>{\raggedright\arraybackslash}p{0.22\linewidth}>{\raggedright\arraybackslash}p{0.14\linewidth}>{\raggedright\arraybackslash}p{0.18\linewidth}>{\raggedright\arraybackslash}p{0.23\linewidth}@{}}
\toprule
\textbf{Date and feast} & \textbf{Pre-1955} & \textbf{1962} & \textbf{Postconciliar} & \textbf{Note}\\
\midrule
\endhead
\bottomrule
\endfoot
24 March: \latin{S.~Gabrielis Archangeli} & \latin{Duplex majus} (\emph{MR}~1920, kalendarium; Mame 1922, pp.~552--554) & III classis (p.~493) & Honoured on 29 September & The \emph{Raccolta} of 1910 lists ``St Gabriel the Archangel, March 18'' among the feasts of the chaplet's indulgences (n.~291)\\
\addlinespace
8 May: \latin{In Apparitione S.~Michaelis Archangeli} & \latin{Duplex majus} (\emph{MR}~1920, kalendarium and proper); Greater Double (Bute, II, p.~867) & Among the Masses \latin{pro aliquibus locis}: \latin{Missa Benedicite}, as on 29 September (p.~[161]) & No celebration of the Archangel on this day & Office lessons of Monte Gargano under Gelasius~I (\emph{BR}, \emph{Pars verna}, pp.~571--572)\\
\addlinespace
29 September: \latin{In Dedicatione S.~Michaelis Archangeli} & \latin{Duplex I classis} (Mame 1922, kalendarium and p.~725); \latin{Duplex 2.~classis} (\emph{BR} 1898, p.~378); Double of the Second Class (Bute, IV, p.~592) & I classis (p.~670) & Saints Michael, Gabriel, and Raphael, Archangels: Feast (\emph{OLM} n.~647) & The Leonine Sacramentary keeps on 30 September the \latin{natale basilicae Angeli in Salaria}; the Gregorian, the \latin{dedicatio basilicae sancti Angeli} on 29 September\\
\addlinespace
2 October: \latin{Ss.~Angelorum Custodum} & \latin{Duplex majus} (Mame 1922, kalendarium and pp.~726--729; \emph{BR} 1898, p.~397) & III classis (p.~672) & The Holy Guardian Angels: Memorial (\emph{OLM} n.~650) & Permitted by Paul~V (1608) and made of precept on this day by Clement~X (1670), as \emph{The Liturgical Year} relates\\
\addlinespace
24 October: \latin{S.~Raphaelis Archangeli} & \latin{Duplex majus} (Mame 1922, kalendarium and pp.~741--743); Greater Double in Bute's Breviary for England (IV, kalendar) & III classis (p.~695) & Honoured on 29 September & The \emph{Raccolta} lists the day among the feasts of the chaplet's indulgences (n.~291)\\
\end{longtable}
\end{center}

\subsection*{The Proper Texts of the Feasts}

In the pre-1955 and 1962 Missals every feast of the angels opens with the
Introit \latin{Benedicite Dominum, omnes Angeli ejus} (Ps 102[103]:20) and
closes with the Communion \latin{Benedicite, omnes Angeli Domini, Dominum}
(Dan 3:58). The proper orations and readings are these (\emph{MR}~1920,
Mame 1922; \emph{MR}~1962 as indexed).

\begin{fourstable}{Feast}{Collect}{Lesson and Gospel}{Proper chants}
29 September (and 8 May) & \latin{Deus, qui miro ordine Angelorum
ministeria hominumque dispensas} & Apoc 1:1--5; Matt 18:1--10 &
Alleluia \latin{Sancte Michael Archangele, defende nos in praelio};
in paschal time \latin{Concussum est mare}; Offertory \latin{Stetit
Angelus} (Apoc 8:3--4)\\
2 October & \latin{Deus, qui ineffabili providentia sanctos Angelos tuos
ad nostram custodiam mittere dignaris} & Exod 23:20--23; Matt 18:1--10 &
Gradual \latin{Angelis suis Deus mandavit de te} (Ps 90[91]:11--12);
Offertory \latin{Benedicite Dominum} (Ps 102[103]:20--21)\\
24 March & \latin{Deus, qui inter ceteros Angelos, ad annuntiandum
incarnationis tuae mysterium, Gabrielem Archangelum elegisti} & Dan
9:21--26; Luke 1:26--38 & Tract \latin{Ave, Maria} (Luke 1:28, 42, 31,
35); Offertory \latin{Stetit Angelus}\\
24 October & \latin{Deus, qui beatum Raphaelem Archangelum Tobiae famulo
tuo comitem dedisti in via} & Tob 12:7--15; John 5:1--4 & Gradual
\latin{Angelus Domini Raphael apprehendit et ligavit daemonem} (Tob 8:3);
Offertory \latin{Stetit Angelus}\\
\end{fourstable}

\subsection*{The Postconciliar Readings}

\begin{fourstable}{Celebration}{First reading}{Psalm and acclamation}{Gospel}
29 September, Saints Michael, Gabriel, and Raphael (Feast) & Dan
7:9--10, 13--14, or Apoc 12:7--12a & Ps 137[138]:1--5; Ps 102[103]:21 &
John 1:47--51\\
2 October, The Holy Guardian Angels (Memorial) & Exod 23:20--23 & Ps
90[91]:1--6, 10--11; Ps 102[103]:21 & Matt 18:1--5, 10\\
\end{fourstable}

The Directory of 2001 names the two postconciliar days and the votive Mass
of the Holy Angels among the Church's celebrations of the angels
(\emph{Directory on Popular Piety} 215).
